\begin{ex}
Tìm số trung bình cộng của 20 và 30.
\shortans{25}
\loigiai{
$(20 + 30) : 2 = 25$.
}
\end{ex}

\begin{ex}
Tìm số trung bình cộng của 15 và 25.
\shortans{20}
\loigiai{
$(15 + 25) : 2 = 20$.
}
\end{ex}

\begin{ex}
Tìm số trung bình cộng của 10, 20 và 30.
\shortans{20}
\loigiai{
$(10 + 20 + 30) : 3 = 20$.
}
\end{ex}

\begin{ex}
Tìm số trung bình cộng của 12, 14 và 16.
\shortans{14}
\loigiai{
$(12 + 14 + 16) : 3 = 14$.
}
\end{ex}

\begin{ex}
Tìm số trung bình cộng của 25, 35 và 45.
\shortans{35}
\loigiai{
$(25 + 35 + 45) : 3 = 35$.
}
\end{ex}

\begin{ex}
Tìm số trung bình cộng của 4 số: 10, 20, 30, 40.
\shortans{25}
\loigiai{
$(10 + 20 + 30 + 40) : 4 = 100 : 4 = 25$.
}
\end{ex}

\begin{ex}
Tìm số trung bình cộng của 4 số: 15, 25, 35, 45.
\shortans{30}
\loigiai{
$(15 + 25 + 35 + 45) : 4 = 120 : 4 = 30$.
}
\end{ex}

\begin{ex}
Tìm số trung bình cộng của 5 số: 1, 2, 3, 4, 5.
\shortans{3}
\loigiai{
$(1 + 2 + 3 + 4 + 5) : 5 = 15 : 5 = 3$.
}
\end{ex}

\begin{ex}
Tìm số trung bình cộng của 5 số: 10, 20, 30, 40, 50.
\shortans{30}
\loigiai{
$(10 + 20 + 30 + 40 + 50) : 5 = 150 : 5 = 30$.
}
\end{ex}

\begin{ex}
Bạn Mai có 12 bông hoa, bạn Lan có 18 bông hoa. Trung bình mỗi bạn có bao nhiêu bông hoa?
\shortans{15}
\loigiai{
$(12 + 18) : 2 = 15$ bông.
}
\end{ex}

\begin{ex}
Hùng cân nặng 32 kg, Dũng cân nặng 36 kg, Nam cân nặng 40 kg. Trung bình mỗi bạn cân nặng bao nhiêu ki-lô-gam?
\shortans{36}
\loigiai{
$(32 + 36 + 40) : 3 = 36$ kg.
}
\end{ex}

\begin{ex}
Bác An ngày thứ nhất thu hoạch được 45 kg thóc, ngày thứ hai thu hoạch được 55 kg thóc. Trung bình mỗi ngày bác thu được bao nhiêu ki-lô-gam thóc?
\shortans{50}
\loigiai{
$(45 + 55) : 2 = 50$ kg.
}
\end{ex}

\begin{ex}
Một đội xe chở hàng: xe 1 chở 35 tạ, xe 2 chở 40 tạ, xe 3 chở 45 tạ. Trung bình mỗi xe chở bao nhiêu tạ?
\shortans{40}
\loigiai{
$(35 + 40 + 45) : 3 = 40$ tạ.
}
\end{ex}

\begin{ex}
Khối 4 có 3 lớp: lớp 4A có 30 học sinh, lớp 4B có 32 học sinh, lớp 4C có 34 học sinh. Trung bình mỗi lớp có bao nhiêu học sinh?
\shortans{32}
\loigiai{
$(30 + 32 + 34) : 3 = 32$ học sinh.
}
\end{ex}

\begin{ex}
Trung bình cộng của hai số là 45. Tổng của hai số đó là bao nhiêu?
\shortans{90}
\loigiai{
$45 \times 2 = 90$.
}
\end{ex}

\begin{ex}
Trung bình cộng của ba số là 20. Tổng của ba số đó là bao nhiêu?
\shortans{60}
\loigiai{
$20 \times 3 = 60$.
}
\end{ex}

\begin{ex}
Trung bình cộng của bốn số là 25. Tổng của bốn số đó là bao nhiêu?
\shortans{100}
\loigiai{
$25 \times 4 = 100$.
}
\end{ex}

\begin{ex}
Trung bình cộng của hai số là 50, biết một số là 40. Số còn lại là bao nhiêu?
\shortans{60}
\loigiai{
Tổng = $50 \times 2 = 100$. Số kia = $100 - 40 = 60$.
}
\end{ex}

\begin{ex}
Trung bình cộng của hai số là 35, biết một số là 25. Số còn lại là bao nhiêu?
\shortans{45}
\loigiai{
Tổng = 70. Số kia = $70 - 25 = 45$.
}
\end{ex}

\begin{ex}
Trung bình cộng của ba số là 30. Biết tổng của hai số đầu là 50. Số thứ ba là bao nhiêu?
\shortans{40}
\loigiai{
Tổng cả 3 số = $30 \times 3 = 90$. Số thứ 3 = $90 - 50 = 40$.
}
\end{ex}

\begin{ex}
Số trung bình cộng của 5 số lẻ liên tiếp bắt đầu từ 1 là bao nhiêu? (1, 3, 5, 7, 9).
\shortans{5}
\loigiai{
Số ở chính giữa là 5.
}
\end{ex}

\begin{ex}
Số trung bình cộng của 5 số chẵn liên tiếp bắt đầu từ 2 là bao nhiêu? (2, 4, 6, 8, 10).
\shortans{6}
\loigiai{
Số ở chính giữa là 6.
}
\end{ex}

\begin{ex}
Điểm kiểm tra 4 môn của bạn An là: Toán 9, Tiếng Việt 8, Khoa học 9, Lịch sử 10. Điểm trung bình của An là bao nhiêu? (Tính tròn số nguyên).
\shortans{9}
\loigiai{
$(9 + 8 + 9 + 10) : 4 = 36 : 4 = 9$ điểm.
}
\end{ex}

\begin{ex}
Trung bình cộng của hai số là 80, hiệu của hai số là 20. Số lớn là bao nhiêu?
\shortans{90}
\loigiai{
Tổng = 160. Số lớn = $(160 + 20) : 2 = 90$.
}
\end{ex}

\begin{ex}
Số bé trong bài toán trên là bao nhiêu?
\shortans{70}
\loigiai{
$(160 - 20) : 2 = 70$.
}
\end{ex}
